\section{The joint electron--lattice response}\label{app:D.33}

We derive the properties of the joint response \cref{eq:12.1-joint}, the Migdal entry \cref{eq:12.1-Skk}, and its reductions to the four published forms of \cref{tab:12-truncations}. The notation is that of \cref{sec:8.6,app:D.28}: adiabatic electronic states $\Psi_j(u)$ with energies $E_j(u)$, $\Delta E_j = E_j - E_0$, the derivative coupling $\vect{d}^I_{j0} = \langle\Psi_j|\nabla_{\vect{X}_I}\Psi_0\rangle$ and the force matrix element $\vect{F}^I_{j0} = \langle\Psi_j|\nabla_{\vect{X}_I}V_{eN}|\Psi_0\rangle$.

\subsection{Labels, marginal and positivity}\label{app:D.33-labels}

Let $H_{\mathrm{el}} = \sum_\Lambda E^{\mathrm{el}}_\Lambda P_\Lambda$ be the operator that multiplies each final-state class $\Lambda$ of \cref{sec:8.7-classes} by its electronic energy, and $H_{\mathrm{lat}} = H - H_{\mathrm{el}}$. Where the classes are good labels, $[P_\Lambda, H] = 0$, and then $H_{\mathrm{el}}$ and $H_{\mathrm{lat}}$ commute with $H$ and with each other. Their joint eigenstates are the eigenstates $|f\rangle$ of $H$, with eigenvalues $E^{\mathrm{el}}_f$ and $E_f - E^{\mathrm{el}}_f$, and \cref{eq:12.1-joint} can be written without reference to a basis,
\begin{equation*}
  S_{ab}(\vect{q};\omega_l,\omega_e) = \frac{2\pi}{\Vol}\sum_ip_i\,\bigl\langle i\big|O_a^\dagger\,\delta\bigl(\omega_e + E^{\mathrm{el}}_i - H_{\mathrm{el}}\bigr)\,\delta\bigl(\omega_l + E_i - E^{\mathrm{el}}_i - H_{\mathrm{lat}}\bigr)\,O_b\big|i\bigr\rangle .
\end{equation*}
Integrating over $\omega_e$ at fixed $\omega = \omega_l + \omega_e$ turns the product of the two delta functions into $\delta(\omega + E_i - H_{\mathrm{el}} - H_{\mathrm{lat}}) = \delta(\omega + E_i - H)$, since the two operators commute, and leaves \cref{eq:4.2-Sab}; this is \cref{eq:12.1-marginal}. For any complex vector $c_a$, the contraction $\sum_{ab}c_a^*S_{ab}c_b$ at a point of the plane is $(2\pi/\Vol)\sum_{i,f}p_i|\langle f|\sum_bc_bO_b|i\rangle|^2$ times the two delta functions, a sum of nonnegative terms. The matrix is therefore positive semidefinite at every $(\omega_l, \omega_e)$, and the bound $|S_{e\kappa}|^2 \le S_{ee}S_{\kappa\kappa}$ of \cref{sec:12.1-def} is its two-by-two minor. Where the labels commute with $H$ only to the accuracy of \cref{eq:8.6-leak} or of the off-shell window, these statements hold on bins of the plane, with the remainders of \cref{eq:8.6-decomp,eq:8.6-bin}.

At $T = 0$ the initial state is the vibrational ground state of the ground surface, with $E^{\mathrm{el}}_i = 0$. A final state in $\mathcal{P}_{\mathrm{ad}}$ has $E^{\mathrm{el}}_f = 0$, which places the electron-following sector on the line $\omega_e = 0$. A final state in $\mathcal{P}_0$ has the zero-phonon energy of its excitation as its total excitation energy, $E_f - E_i = E^{\mathrm{el}}_f$, which places it on the line $\omega_l = 0$. Every other final state has $\omega_e > 0$ and $\omega_l > 0$. This is \cref{eq:12.1-sectors}.

\subsection{The force operator and the amplitude}\label{app:D.33-amp}

Let the electrons feel nucleus $I$, of species $\kappa$, through a potential that moves rigidly with it, $V_{eN} \ni \sum_ev_\kappa(\vect{r}_e - \vect{X}_I)$. With the transform \cref{eq:A.3-space}, $v_\kappa(\vect{r}) = \int\dd^3p\,(2\pi)^{-3}\ee^{-\ii\vect{p}\cdot\vect{r}}v_\kappa(\vect{p})$, and with $v_\kappa(-\vect{p}) = v_\kappa(\vect{p})$ for a spherical potential,
\begin{equation*}
  \sum_ev_\kappa(\vect{r}_e - \vect{X}_I) = \int\frac{\dd^3p}{(2\pi)^3}\;v_\kappa(\vect{p})\,\ee^{-\ii\vect{p}\cdot\vect{X}_I}\sum_e\ee^{\ii\vect{p}\cdot\vect{r}_e}
  = \int\frac{\dd^3p}{(2\pi)^3}\;v_\kappa(\vect{p})\,\ee^{-\ii\vect{p}\cdot\vect{X}_I}\,n_e(\vect{p}) ,
\end{equation*}
after the substitution $\vect{p} \to -\vect{p}$, with $n_e(\vect{p})$ of \cref{eq:A.3-ne}. Only this term of $V_{eN}$ depends on $\vect{X}_I$. Its gradient gives \cref{eq:12.1-force}, and its matrix element between adiabatic states gives
\begin{equation}
  \vect{q}\cdot\vect{F}^I_{j0} = -\ii\int\frac{\dd^3p}{(2\pi)^3}\;(\vect{q}\cdot\vect{p})\,v_\kappa(\vect{p})\,\bigl\langle\Psi_j(u)\big|\,n_e(\vect{p})\,\big|\Psi_0(u)\bigr\rangle\,\ee^{-\ii\vect{p}\cdot\vect{X}_I} .
  \label{eq:D.33-F}
\end{equation}
Both $\ee^{-\ii\vect{p}\cdot\vect{X}_I}$ and the transition density depend on the nuclear positions, and the product is a function of $u$ that multiplies the nuclear wavefunction.

The boost amplitude \cref{eq:8.6-boost} is $-\ii\,\vect{q}\cdot\vect{d}^I_{j0}/(M_\kappa\Delta E_j)$ times $\ee^{\ii\vect{q}\cdot\vect{X}_I}$ between the vibrational states. By \cref{eq:5.1-dcoup}, $\vect{d}^I_{j0} = \vect{F}^I_{j0}/(E_0 - E_j) = -\vect{F}^I_{j0}/\Delta E_j$, so the electronic factor is $\ii\,\vect{q}\cdot\vect{F}^I_{j0}/(M_\kappa\Delta E_j^2)$. The derivation of \cref{eq:D.28-boost} treats the derivative coupling as a constant over the vibrational wavefunctions. Keeping it inside the nuclear matrix element, as a function of $u$, gives the same leading term together with the Herzberg--Teller corrections. With \cref{eq:D.33-F} and $\ii\cdot(-\ii) = 1$,
\begin{equation*}
  A_{fj} = \sum_{I\in\kappa}\int\frac{\dd^3p}{(2\pi)^3}\;\frac{\vect{q}\cdot\vect{p}}{M_\kappa\Delta E_j^2}\,v_\kappa(\vect{p})\,
  \bigl\langle\Psi^{\mathrm{nuc}\prime}\big|\,\langle\Psi_j(u)|n_e(\vect{p})|\Psi_0(u)\rangle\,\ee^{\ii(\vect{q} - \vect{p})\cdot\vect{X}_I}\big|\Psi^{\mathrm{nuc}}\bigr\rangle .
\end{equation*}
The Condon approximation in the frame of the crystal replaces the transition density by its value at $u = 0$, which leaves the operator $\ee^{\ii(\vect{q}-\vect{p})\cdot\vect{X}_I}$ between the vibrational states and gives \cref{eq:12.1-amp}. Its first correction is the gradient of the transition density with respect to the displacements, the Herzberg--Teller term of \cref{eq:8.6-electronic}.

If instead the electrons move rigidly with nucleus $I$, as in an isolated atom, the transition density at nuclear position $\vect{X}_I$ is that at the origin times $\ee^{\ii\vect{p}\cdot\vect{X}_I}$, since translating every electron by $\vect{X}_I$ multiplies $n_e(\vect{p})$ by this phase. The phase cancels the factor $\ee^{-\ii\vect{p}\cdot\vect{X}_I}$ of \cref{eq:D.33-F}, and the operator between the vibrational states becomes $\ee^{\ii\vect{q}\cdot\vect{X}_I}$, a recoil with the full momentum $\vect{q}$. The two cases differ by the dependence of the transition density on $u$ that the frame of the crystal neglects and the frame of the nucleus overstates. Expanded in $\vect{p}\cdot\vect{u}_I$, the two lattice factors agree at zeroth order and differ at first order: the lattice keeps $\vect{q} - \vect{p}$ in one frame and $\vect{q}$ in the other, which matters for the lattice spectrum when $p$ is not small against $q$ (\cref{sec:12.4}).

\subsection{Squared amplitude and crystal momentum}\label{app:D.33-square}

In a box of volume $\Vol$, $\int\dd^3p\,(2\pi)^{-3} \to \Vol^{-1}\sum_{\vect{p}}$. The squared amplitude summed over the nuclei of species $\kappa$ is
\begin{equation*}
  \begin{split}
  |A_{fj}|^2 = \frac{1}{M_\kappa^2\Delta E_j^4\Vol^2}\sum_{\vect{p}\vect{p}'}\;&(\vect{q}\cdot\vect{p})(\vect{q}\cdot\vect{p}')\,v_\kappa(\vect{p})^*v_\kappa(\vect{p}')\,
  \langle\Psi_0|n_e(\vect{p})^\dagger|\Psi_j\rangle\langle\Psi_j|n_e(\vect{p}')|\Psi_0\rangle \\
  &\times\langle\Psi^{\mathrm{nuc}}|n_\kappa(\vect{q}-\vect{p})^\dagger|\Psi^{\mathrm{nuc}\prime}\rangle\langle\Psi^{\mathrm{nuc}\prime}|n_\kappa(\vect{q}-\vect{p}')|\Psi^{\mathrm{nuc}}\rangle ,
  \end{split}
\end{equation*}
with $n_\kappa(\vect{Q}) = \sum_{I\in\kappa}\ee^{\ii\vect{Q}\cdot\vect{X}_I}$, the transition densities at $u = 0$, and the real factor $\vect{q}\cdot\vect{p}$. The final state has $\omega_e = \Delta E_j$ and $\omega_l$ equal to the energy of $\Psi^{\mathrm{nuc}\prime}$ above $\Psi^{\mathrm{nuc}}$. Multiplying by $(2\pi/\Vol)p_i\,\delta(\omega_e - \Delta E_j)\,\delta(\omega_l - E^{\mathrm{nuc}\prime} + E^{\mathrm{nuc}})$ and summing over $j \neq 0$, over $\Psi^{\mathrm{nuc}\prime}$ and over the thermal initial states, the electronic sum is $(\Vol/2\pi)$ times the density response between $n_e(\vect{p})$ and $n_e(\vect{p}')$ at clamped nuclei, and the nuclear sum is $(\Vol/2\pi)$ times the electron-following response between $n_\kappa(\vect{q}-\vect{p})$ and $n_\kappa(\vect{q}-\vect{p}')$. In a crystal the clamped electrons conserve crystal momentum, and their response vanishes unless $\vect{p}$ and $\vect{p}'$ differ by a reciprocal-lattice vector. Writing $\vect{p} = \vect{k} + \vect{G}$ and $\vect{p}' = \vect{k} + \vect{G}'$ with $\vect{k}$ in the zone, the double sum becomes $\sum_{\vect{k}}\sum_{\vect{G}\vect{G}'}$, and $\Vol^{-1}\sum_{\vect{k}} \to \int_{\mathrm{BZ}}\dd^3k\,(2\pi)^{-3}$. Collecting the factors, $(2\pi/\Vol)(\Vol/2\pi)^2\Vol^{-2}\cdot \Vol = (2\pi)^{-1}$, and with $\Delta E_j = \omega_e$ under the delta function we obtain \cref{eq:12.1-Skk}. The units check: the joint response has the units of $S_{ab}(\vect{q},\omega)$ per unit energy.

Of the approximations of \cref{sec:12.1-boost}, the ground-surface final states, which make the lattice factor the electron-following response and $\Delta E_j$ the vertical energy, can be lifted by projecting on the vibrational states of surface $j$, which \cref{we:12-oscillator} does for one oscillator. For a pair delocalized over $N$ cells the displacement of the excited surface is of order $N^{-1}$ per nucleus (\cref{sec:8.6}), and the two choices agree.

\subsection{Diagonal and incoherent reduction}\label{app:D.33-soft}

Keep only $\vect{G} = \vect{G}'$ and the self part of the lattice factor. The self part is exact for the phonon-emitting terms, which are incoherent, and, averaged over the zone, for the zero-phonon term at $q \gtrsim 2q_{\mathrm{BZ}}$. At $q \ll q_{\mathrm{BZ}}$ the coherent zero-phonon part keeps the cross-site terms of the lattice factor and the full matrix $\vect{G} \neq \vect{G}'$, and obeys \cref{eq:8.3-boostward} instead (\cref{sec:12.4}). For one nucleus per site the self response is $S^{\mathrm{ad}}_{\kappa;\vect{G}\vect{G}}(\vect{q} - \vect{k},\omega_l) = 2\pi\nd_\kappa\,s^{\mathrm{inc}}_\kappa(\vect{q} - \vect{k} - \vect{G},\omega_l)$, with $s^{\mathrm{inc}}_\kappa(\vect{Q},\omega)$ the spectrum of one nucleus normalized to $\int\dd\omega\,s^{\mathrm{inc}}_\kappa = 1$, zero-phonon line included, as in \cref{eq:11.3-inc}. At $T = 0$ the diagonal electronic factor is fixed by the fluctuation--dissipation theorem of \cref{sec:10.2}, $S^{ee}_{\vect{G}\vect{G}}(\vect{k},\omega) = (2|\vect{p}|^2/e^2)\,L(\vect{p},\omega)$, with $\vect{p} = \vect{k} + \vect{G}$ and $L(\vect{p},\omega) = \Imag[-(\varepsilon^{-1})_{\vect{G}\vect{G}}(\vect{k},\omega)]$ the loss function at the unreduced momentum. For a vertex $v_\kappa(\vect{p}) = -Z_\kappa e^2/p^2$, the sum over $\vect{G}$ and the integral over the zone combine into an integral over all $\vect{p}$, and dividing by the weight $2\pi\nd_\kappa$ of one nucleus,
\begin{equation}
  \frac{\dd^2P}{\dd\omega_l\,\dd\omega_e} = 4\alpha\,Z_\kappa^2\int\frac{\dd^3p}{(2\pi)^3}\;\frac{(\vect{v}_N\cdot\vect{p})^2}{\omega_e^4\,p^2}\,L(\vect{p},\omega_e)\,s^{\mathrm{inc}}_\kappa(\vect{q} - \vect{p},\omega_l) ,
  \label{eq:D.33-diag}
\end{equation}
with $\vect{v}_N = \vect{q}/M_\kappa$ and $e^2 = 4\pi\alpha$: the factors are $(2\pi)^{-1}\cdot Z_\kappa^2e^4p^{-4}\cdot2p^2e^{-2} = Z_\kappa^2e^2/(\pi p^2)$. \Cref{fig:12-plane} evaluates \cref{eq:D.33-diag} for silicon with $Z_\kappa = 4$. Let the momentum taken by the electrons be small against the kick, $p \ll q$. The lattice factor then depends on $\vect{q}$ alone, its integral over $\omega_l$ is one, and
\begin{equation}
  \frac{\dd P}{\dd\omega_e} = 4\alpha\int\frac{\dd^3p}{(2\pi)^3}\;Z_\kappa^2\,\frac{(\vect{v}_N\cdot\vect{p})^2}{\omega_e^4\,p^2}\,L(\vect{p},\omega_e) .
  \label{eq:D.33-4alpha}
\end{equation}
With the charge $Z_\kappa$ replaced by the momentum-dependent ion charge $Z_{\mathrm{ion}}(p)$ of the nucleus and its frozen core, this is the soft-limit formula of Knapen, Kozaczuk and Lin~\cite[Eq.~(6)]{knapen2021migdal}, coefficient $4\alpha$ included. Their full expression~\cite[Eq.~(3)]{knapen2021migdal} keeps the impulse lattice factor at $\vect{q} - \vect{p}$, a free recoil smeared by the motion of the nucleus in its well, which is \cref{eq:D.33-diag} with the impulse form of $s^{\mathrm{inc}}_\kappa$ (\cref{sec:11.4}). Its energy bracket carries the final momentum of the nucleus in the numerator and drops $p^2/2M_\kappa$ in both denominators, which is consistent only in the soft limit; the exact free-recoil bracket, with $p^2/2M_\kappa$ in both denominators and the numerator $\vect{q}\cdot\vect{p}/M_\kappa$, is \cref{eq:12.3-bracket}.

\subsection{The isolated atom}\label{app:D.33-atom}

For an atom the electrons move with the nucleus, and by \cref{app:D.33-amp} the lattice factor is $\ee^{\ii\vect{q}\cdot\vect{X}}$. The electronic factor $\vect{q}\cdot\vect{F}_{j0}/(M\Delta E_j^2)$ is then evaluated with the atom at the origin, where the adiabatic states depend on $\vect{X}$ only through $\vect{r}_e - \vect{X}$. \Cref{sec:6.1} showed that this makes the derivative coupling a dipole matrix element, $\vect{d}_{j0} = m_e\Delta E_j\langle\Psi_j|\sum_e\vect{r}_e|\Psi_0\rangle$. The amplitude becomes $-\ii m_e\vect{v}_N\cdot\langle\Psi_j|\sum_e\vect{r}_e|\Psi_0\rangle$, the first-order term of Migdal's boost \cref{eq:6.1-boost}. The boost itself is exact for a free atom in the sudden limit: the electrons are left in $\Psi_0$ while the nucleus moves with $\vect{v}_N$, and the probability of the excited state $j$ is $|\langle\Psi_j|\exp(-\ii m_e\vect{v}_N\cdot\sum_e\vect{r}_e)|\Psi_0\rangle|^2$ to all orders in $v_N$. This is the electron-cloud transition factor of Ibe and collaborators~\cite[Eqs.~(55) and (57)]{ibe2018migdal}, whose argument is $m_e$ times the final velocity of the atom~\cite[Eq.~(58)]{ibe2018migdal}, $\vect{v}_N$ up to corrections of order $m_e/M_\kappa$. The lattice factor is the free recoil of the atom. Its mass is that of the atom, not of the bare nucleus, since the drag mass of an atom is $\delta M = N_em_e$ (\cref{sec:8.6}); this is the boost mass $m_A$ of \cref{sec:12.2}. With the free recoil the energy is $\omega_l = q^2/2m_A$ for every Migdal event, and the joint response is a line in the plane.

\subsection{Bound nuclei and the virtual phonon}\label{app:D.33-phonon}

Below the impulse onset the lattice factor of \cref{eq:12.1-Skk} is the harmonic multiphonon response at $\vect{q} - \vect{k} - \vect{G}$, \cref{eq:11.3-exact} or its incoherent form \cref{eq:11.3-inc}. With the diagonal of the loss function and the vertex of a pseudo-ion, \cref{eq:12.1-Skk} then has the form of the rate of Berghaus and collaborators~\cite[Eq.~(3.8)]{berghaus2023migdal}, whose effective interaction $\nabla_IH_{\chi L}\cdot\nabla_IH_{eL}/(m_N\omega^2)$~\cite[Eq.~(3.6)]{berghaus2023migdal} is the product of the kick and the force \cref{eq:12.1-force} divided by $M_\kappa\omega_e^2$. Their lattice structure factor is evaluated at $\vect{q} - \vect{k}_e - \vect{K}$, the frame of the crystal of \cref{app:D.33-amp}.

The energy denominator $M_\kappa\omega_e^2$ assumes a nucleus that is free during the electronic time $1/\omega_e$. A nucleus bound in the lattice is pulled back by its neighbors, and \cref{sec:6.1} found that the sudden amplitude of a nucleus kicked in a harmonic well of frequency $\omega_{\mathrm{ph}}$ acquires the factor $[1 - (\omega_{\mathrm{ph}}/\omega_e)^2]^{-1}$. Resolved into phonon branches, the inertial response $1/(M_\kappa\omega_e^2)$ of the kicked nucleus becomes the phonon propagator
\begin{equation}
  \frac{\vect{q}\cdot\vect{p}}{M_\kappa\omega_e^2} \;\to\; \frac{1}{N}\sum_{\vect{Q}\nu}\frac{(\vect{q}\cdot\vect{e}_{\kappa\nu}(\vect{Q}))\,(\vect{p}\cdot\vect{e}_{\kappa\nu}(\vect{Q})^*)}{M_\kappa\,(\omega_e^2 - \omega_{\vect{Q}\nu}^2)}
  = \frac{\vect{q}\cdot\vect{p}}{M_\kappa\omega_e^2}\Bigl[1 + O\Bigl(\frac{\omega_{\mathrm{ph}}^2}{\omega_e^2}\Bigr)\Bigr] ,
  \label{eq:D.33-virtual}
\end{equation}
where the equality uses the completeness of the polarization vectors of each $\vect{Q}$, $\sum_\nu\vect{e}_{\kappa\nu}\vect{e}_{\kappa\nu}^* = \mathbf{1}$, on the leading term. This is the virtual phonon of Liang and collaborators~\cite[Eqs.~(3.1) and (3.2)]{liang2022phonon}, which they reduce to $1/(\varepsilon_i - \varepsilon_j)^2$ for gaps large against the phonon energies, and whose multiphonon sum~\cite[Eqs.~(3.3) and (3.4)]{liang2022phonon} is the incoherent lattice factor with its zero-phonon term. Averaged over the phonon density of states, the correction in the probability is $2\langle\omega_{\mathrm{ph}}^2\rangle/\omega_e^2$, which in silicon at the gap, with $\langle\omega_{\mathrm{ph}}^2\rangle^{1/2} = \SI{44.7}{meV}$, is $2.9\times10^{-3}$, of the order of the Condon and screening corrections that the same formulas neglect. \Cref{sec:12.4} shows that for a harmonic lattice the propagator form is exact at first order in $\vect{p}$, treats it below the impulse onset and splits it into its zero-phonon and phonon-emitting terms.
